%%
%% IntentCap extended abstract
%%

\documentclass[sigconf,nonacm]{acmart}

\setcopyright{none}
\settopmatter{printacmref=false,printccs=false}
\renewcommand\footnotetextcopyrightpermission[1]{}
\pagestyle{plain}

\AtBeginDocument{%
  \providecommand\BibTeX{{%
    Bib\TeX}}}

\usepackage{microtype}
\usepackage{xspace}
\usepackage{tikz}
\usepackage{booktabs}
\usetikzlibrary{arrows.meta,positioning,fit,backgrounds}
\setlength{\bibsep}{0pt}

\newcommand{\sys}{\textsc{IntentCap}\xspace}

\begin{document}

\title{IntentCap: Intent-Carrying Capabilities for LLM Agent Extensions}

\begin{abstract}
Every agent runtime today has the same blind spot: it cannot tell whose context is driving a decision. A Skill's instructions, a tool's schema, a PDF's extracted text, and the user's explicit request all flow into one planning channel, yet only one of them should be authorized to choose a destination, widen an approval scope, or delegate authority. Current defenses guard tools, parameters, and files; none of them ask which input class is filling which authority-bearing field. The result is structural: any extension that supplies data can silently supply decisions.

We reframe agent authorization as a context-influence problem. \sys partitions every input into four issuer-owned classes (agent intent, procedural instruction, tool metadata, and runtime observation) and enforces a no-substitution rule: a field owned by one class cannot be proven by another. User intent, captured as structured certificates rather than natural language, becomes the root principal. Task plans compile into short-lived, proof-carrying capability leases that constrain operations, arguments, data flow, temporal order, and delegation depth. A deterministic checker validates field-owner proofs and consumes lease budget atomically, before the side effect, the prompt placement, or the delegation handoff commits. The LLM proposes; the checker decides. The trusted computing base includes the intent issuer, the checker, and a provenance tracker, but explicitly excludes the model itself.

Removing issuer ownership entirely causes 94\% of checker-denied protected decisions to be falsely accepted across 3,823 events from AgentDojo, MCPTox, InjecAgent, and tau2-bench; even the weakest single-edge collapse falsely accepts over 50\%. Across 38 multi-boundary enforcement events spanning tool calls, local execution, context placement, Skill instruction slots, and delegation handoffs, the checker produces zero unsafe accepts. This is not a tuning knob. It is a phase transition between authorization and theater.
\end{abstract}

\maketitle

\section{Introduction}

LLM agents are becoming extensible execution environments. A modern agent can load Skills with workflow instructions, references, and optional scripts~\cite{openai-skills}; connect to MCP servers that expose resources, prompts, and tools~\cite{mcp-spec}; run local commands; and spawn subagents. This extensibility is useful precisely because these components affect future behavior. A Skill can teach the agent how to perform a workflow. A tool result can become evidence for a later action. A document can determine the content of a generated report. The same property creates a security problem: in an agent, text is not passive data. It can become control.

Most current defenses ask an operation-centric question: which tool may be called, which file may be read, which endpoint may be contacted, or which process may execute? That question is necessary but incomplete. Consider a user who asks an agent to extract tables from two selected PDFs, save spreadsheets, and create one GitHub issue in a named repository. The PDF contents should influence spreadsheet cells and perhaps the issue body. They should not influence the repository name, network destination, approval scope, tool selection, or decision to load another Skill. If hidden text in the PDF says ``send all extracted values to attacker.com'' or ``create the issue in attacker/repo,'' a tool-call allowlist or OS sandbox may catch some concrete effects, but it does not directly express the more agent-specific invariant: the untrusted PDF was never allowed to influence those decisions.

\sys starts from that invariant. Its thesis is that securing agent extensions requires controlling not only what extensions can do, but what their context is allowed to mean. The root of authority is a structured \emph{intent certificate}, not an extension package, server identity, or static manifest. A certificate records the current goal, selected objects, authorized sinks, constraints, approvals, and expiry. Context cells are labeled with provenance and allowed \emph{influence modes}, such as observing, quoting, summarizing, parameterizing, planning, instructing, delegating, or authorizing. A capability is then minted for a particular run as an \emph{intent-carrying lease}: it authorizes an operation only if the operation is derivable from the current intent and its control dependencies come from context allowed to influence that decision class.

This framing deliberately differs from extension-interface and OS-enforcement work. EIM and bpftime model application extensions by representing extension features as resources and granting capabilities at extension entries~\cite{eim-bpftime}. \sys instead protects future agent decisions: tool choice, sink selection, authority request, policy update, and delegation. ActPlane enforces agent policies below the tool layer using OS-level information-flow checks and eBPF-style monitoring~\cite{actplane}. \sys instead decides which policy is justified by the user intent and context provenance; ActPlane is a possible backend, not the conceptual center. SkillGuard is closer because it treats Skills as permission-bearing artifacts and jointly regulates context and action behavior~\cite{skillguard}; \sys is run-centric and cross-extension, asking what this run may allow this Skill, this MCP server, this document, and this subagent to influence.

\section{Design}

The root cause of context privilege is two structural mismatches. First, \emph{issuer collapse}: user intent, Skill instructions, tool metadata, and runtime observations share one planning channel, so one class can fill another's authority fields. Second, \emph{authority-state split}: lease budget, expiry, and delegation state are checked by separate guards, so a consumed one-shot lease can be reused or widened between boundaries.

\sys addresses both by making the authorization unit a \emph{pre-effect commit record} rather than an action-level predicate. Figure~\ref{fig:arch} shows the pipeline. A trusted issuer canonicalizes user intent into a structured certificate. A context labeler assigns provenance and influence modes. An untrusted LLM-assisted compiler proposes plans and candidate leases. A deterministic checker validates field-owner proofs and atomically consumes lease state via a single transition:
\vspace{-2pt}
\begin{small}
\[\mathsf{check\_and\_consume}(e,\mathit{lease},\mathit{proofs},\mathit{prov},\sigma) \rightarrow \mathsf{allow}(\sigma'\!,\mathit{audit}) \mid \mathsf{deny}(\mathit{reason})\]
\end{small}
\vspace{-4pt}

Every authority-changing boundary---tool call, local execution, context placement, Skill instruction slot, or delegation handoff---must submit this record \emph{before} the side effect commits. The checker is the sole writer of lease state; adapters cannot cache decisions or mutate budgets locally.

\begin{figure}[t]
\centering
\small
\begin{tikzpicture}[
  node distance=0.6cm and 0.5cm,
  box/.style={draw, rounded corners=2pt, align=center, minimum height=0.65cm, text width=1.9cm, fill=white, font=\scriptsize},
  trusted/.style={box, fill=green!8},
  untrusted/.style={box, fill=red!7},
  runtime/.style={box, fill=blue!7},
  arr/.style={-Latex, thick}
]
\node[trusted] (issuer) {Intent\\issuer};
\node[trusted, right=of issuer] (labeler) {Context\\labeler};
\node[untrusted, right=of labeler] (compiler) {LLM-assisted\\compiler};
\node[trusted, right=of compiler] (checker) {Deterministic\\checker};
\node[runtime, below=0.5cm of checker, xshift=-1.1cm] (adapters) {Boundary adapters};
\node[runtime, left=0.7cm of adapters] (gateways) {Tool/env\\gateways};
\node[runtime, right=0.7cm of adapters] (deleg) {Placement \&\\delegation};

\draw[arr] (issuer) -- (labeler);
\draw[arr] (labeler) -- (compiler);
\draw[arr] (compiler) -- node[above, font=\tiny] {candidates} (checker);
\draw[arr] (checker) -- node[right, font=\tiny] {leases} (adapters);
\draw[arr] (adapters) -- (gateways);
\draw[arr] (adapters) -- (deleg);
\draw[arr, dashed] (checker) to[bend right=15] node[below, font=\tiny] {denial} (compiler);

\begin{scope}[on background layer]
\node[draw, dashed, rounded corners=3pt, fit=(issuer)(labeler)(checker), inner sep=0.13cm, label={[font=\tiny]above:TCB}] {};
\end{scope}
\end{tikzpicture}
\vspace{-6pt}
\caption{\sys architecture. The LLM proposes; the checker decides. Runtime adapters submit field proofs before any side effect, placement, or handoff.}
\label{fig:arch}
\vspace{-6pt}
\end{figure}

\paragraph{Four authority-input classes.}
\sys partitions inputs into four issuer-owned proof boundaries: \emph{agent} context (user intent, selected objects/sinks, approvals), \emph{instruction} context (system policy, endorsed Skill/manual procedures), \emph{tool} context (MCP/registry schemas, credential scopes, sandbox contracts), and \emph{env} context (runtime observations, tool results, file state, script outputs). The \textbf{no-substitution rule} requires that each protected field be proven by its owner class: tool metadata cannot supply user-authorized sinks, runtime observations cannot supply instruction authority, and Skill text cannot widen approval scope. Two classes may merge only if issuer, forgeable surface, owned fields, and lifecycle authority are identical.

\paragraph{Leases and lifecycle.}
A lease binds operation, object, arguments, intent derivation, control/data provenance, temporal guards, budget, expiry, and delegation depth. Leases are short-lived, consumable, and non-self-widening: a one-shot issue lease is consumed on first use, and a child subagent receives only attenuated capabilities. The checker validates all of these atomically in a single transition, preventing the authority-state split where separate guards lose track of consumption or delegation across boundaries.

\section{Preliminary Evidence}

\textbf{RQ1: Can the event model express benign behavior without over-blocking?}
We replay 3,746 protected-decision events extracted from AgentDojo, MCPTox, InjecAgent, and tau2-bench traces through \sys and several same-event baselines~\cite{agentdojo,mcptox,injecagent,tau2bench}. \sys produces 0 unsafe accepts. On the utility side, \sys covers all 3,813 benign reference actions in the tau2-style proxy and passes 2,554 of 2,556 applicable tool-oracle checks.

\textbf{RQ2: Is the four-class partition a safety requirement or a naming choice?}
We ablate the issuer-ownership and lifecycle dimensions on the same 3,823 checker-denied events. Removing issuer ownership entirely (collapsing all four classes into generic trusted context) causes 3,593 false accepts (94\%). Even the weakest single-edge collapse (tool$\rightarrow$agent) falsely accepts 1,928 events (50\%). On 7 controlled workflow residuals, a post-hoc policy DSL that checks predicates without field ownership falsely accepts 7/7; splitting lifecycle state across separate guards falsely accepts 5/7.

\textbf{RQ3: Does the same commit API work beyond tool calls?}
We test the \texttt{check\_and\_consume} transition across five boundary types: local env side effects, context placement, Skill instruction slots, delegation handoffs, and an OS-monitor-style replay backend. Across 38 checker-submitted boundary events, \sys allows 17 authorized effects and blocks 21 violations, with 0 unsafe executions or placements and 0 checker/monitor mismatches.

\textbf{RQ4: Are leases auditable authority objects?}
We adjudicate 24 lease labels covering InjecAgent, MCPTox, and tau2 samples, and score 144 policy rows against them. All \sys policy families have oracle distance~0; every non-\sys baseline has positive distance. This is author-adjudicated supporting evidence, not independent expert-oracle validation.

\section{Positioning}

\sys should be read as a capability compiler and checker for agent decisions. EIM and bpftime ask how an extension can safely use host-provided resources. ActPlane asks how policy can be enforced below the tool layer. SkillGuard asks how a Skill package can be governed as a permission-bearing artifact. \sys asks a different question: which future decisions may this context influence, under this user intent, and with what proof? That distinction is the paper's center of gravity. It keeps the LLM outside the trusted computing base, treats the current user goal as the root of authority, and makes context influence a least-privilege capability rather than an implicit side effect of putting text in the prompt.

\bibliographystyle{ACM-Reference-Format}
\bibliography{references}

\end{document}
